\DocumentMetadata{
  pdfversion=2.0,
  pdfstandard=ua-2,
  lang=en-US,
  tagging=on,
  tagging-setup={math/setup=mathml-SE,tabsorder=structure}
}
\documentclass[11pt,letterpaper]{article}
\usepackage[margin=0.68in,includefoot]{geometry}
\usepackage{newcomputermodern}
\usepackage{amsmath,amscd,mathtools}
\providecommand{\square}{\mdlgwhtsquare}
\usepackage[hidelinks]{hyperref}
\hypersetup{
  pdftitle={Analysis First-Year Exam, January 2018, Part 2},
  pdfauthor={University of Florida Department of Mathematics},
  pdfsubject={Graduate mathematics examination},
  pdfdisplaydoctitle=true,
  bookmarksopen=true
}
\setcounter{secnumdepth}{0}
\setlength{\parindent}{0pt}
\setlength{\parskip}{0.28em}
\setlength{\emergencystretch}{3em}
\setlength{\leftmargini}{2.15em}
\setlength{\leftmarginii}{2.65em}
\setlength{\leftmarginiii}{3em}
\renewcommand{\labelenumi}{\textbf{\theenumi.}}
\pagestyle{empty}
\raggedbottom
\allowdisplaybreaks
\newenvironment{examproblems}[1]{%
  \begin{enumerate}%
  \setcounter{enumi}{#1}%
  \setlength{\topsep}{0.35em}%
  \setlength{\partopsep}{0pt}%
  \setlength{\itemsep}{0.48em}%
  \setlength{\parsep}{0.18em}%
}{\end{enumerate}}
\newenvironment{examsubparts}{%
  \begin{enumerate}%
  \setlength{\topsep}{0.18em}%
  \setlength{\partopsep}{0pt}%
  \setlength{\itemsep}{0.16em}%
  \setlength{\parsep}{0.1em}%
}{\end{enumerate}}
\newenvironment{examinstructions}{%
  \par\begingroup\itshape
}{%
  \par\endgroup\vspace{0.18em}
}
\makeatletter
\renewcommand\section{\@startsection{section}{1}{0pt}%
  {-1.25ex plus -.25ex minus -.1ex}%
  {0.55ex plus .1ex}%
  {\normalfont\large\bfseries}}
\renewcommand{\@maketitle}{%
  \newpage\null\vskip -1.2em%
  {\footnotesize\bfseries
    \href{https://www.ufl.edu/}{University of Florida}, \href{https://math.ufl.edu/}{Department of Mathematics}\par}%
  \vskip 0.18em%
  \tagstructbegin{tag=Title}%
    {\LARGE\bfseries\href{https://gma.math.ufl.edu/exams/analysis-first-year/}{Analysis First-Year Exam}, January 2018, Part 2\par}%
  \tagstructend%
  \vskip 0.35em%
  \tagmcbegin{artifact}%
    \hrule height 1pt%
  \tagmcend%
  \vskip 0.55em%
}
\makeatother
\title{Analysis First-Year Exam, January 2018, Part 2}
\author{}
\date{}
\begin{document}
\maketitle
\thispagestyle{empty}
\begin{examinstructions}
Answer FOUR questions in detail. State carefully any results used without proof.
\end{examinstructions}
\begin{examproblems}{0}
\item
Let \(f:[0,1]\to\mathbb{R}\) be Riemann-integrable. Prove that if the Riemann integral \(\int_0^1 f\) is nonzero then there exist \(\delta>0\) and a nonempty open interval \(I\subseteq[0,1]\) such that \(|f|\geq\delta\) throughout~\(I\).
\item
Let \(f_n(t)=t^n(1-t)^n\) whenever \(0\leq t\leq 1\) and~\(n\) is a positive integer. Does the sequence \((f_n:n>0)\) converge on \([0,1]\) pointwise? Uniformly? Justify.
\item
Let \(f:[0,1]\to\mathbb{R}\) be continuous. Assume that \(\int_0^1 f(t)e^{-nt}\,dt=0\) whenever~\(n\) is a non-negative integer. Does it follow that~\(f\) is identically zero? Does the answer change if the exponent \(-nt\) is replaced by \(+nt\)? By \(2nt\)? Justify.
\item
Given that \((f_n:n\geq 0)\) is a sequence of measurable real-valued functions, explain why each of the following sets is (or is not) measurable:
\begin{examsubparts}
\item[\textnormal{(i)}]
\(X=\{\omega:f_n(\omega)\to\infty\text{ as }n\to\infty\}\);
\item[\textnormal{(ii)}]
\(Y=\{\omega:f_n(\omega)^2<f_n(\omega)\text{ for each }n\geq 0\}\);
\item[\textnormal{(iii)}]
\(Z=\left\{\omega:\sum_{n=0}^{\infty}f_n(\omega)\text{ converges absolutely}\right\}\).
\end{examsubparts}
\item
Let \((f_n:n\geq 1)\) be a sequence of continuous functions from \([0,1]\) to \([0,1]\) and assume that this sequence converges to~\(f\) pointwise. True or false (proof or counterexample)?
\begin{examsubparts}
\item[\textnormal{(i)}]
\(f\) is Riemann-integrable and (Riemann integrals) \(\int_0^1 f_n\to\int_0^1 f\);
\item[\textnormal{(ii)}]
\(f\) is Lebesgue-integrable and (Lebesgue integrals) \(\int_0^1 f_n\to\int_0^1 f\).
\end{examsubparts}
\end{examproblems}
\end{document}
